%% tikzlibrarycommtikz.paths.code.tex
%% Signal Domain Paths, Background Layer Connector Rendering, Mid-Path Arrows,
%% Curved Connectors, Waveguide Bends, and Inline Path Placement.

% ==============================================================================
% 1. Background Layer Hook for Connector Paths
%    Draws signal connector lines on the 'background' layer behind all main
%    component nodes while keeping inline nodes and labels on the 'main' layer.
% ==============================================================================
\newbox\commtikz@savedfigbox
\newif\ifcommtikz@pathonbg
\def\commtikz@mainlayername{main}

\ifx\commtikz@orig@tikz@finish\@undefined
  \let\commtikz@orig@tikz@finish=\tikz@finish
  \def\tikz@finish{%
    \iftikz@is@node
      \commtikz@orig@tikz@finish
    \else
      \pgfsyssoftpath@getcurrentpath\commtikz@temppath
      \pgfprocessround{\commtikz@temppath}{\commtikz@temppath}%
      \pgfsyssoftpath@setcurrentpath\commtikz@temppath
      \pgfsetcornersarced{\pgfpointorigin}%
      \ifcommtikz@pathonbg
        \ifx\pgfonlayer@name\commtikz@mainlayername
          \commtikz@pathonbgfalse
          \global\setbox\commtikz@savedfigbox=\box\tikz@figbox
          \pgfonlayer{background}%
          \pgfsetlinewidth{\tikzscope@linewidth}%
          \begingroup
          \commtikz@orig@tikz@finish
          \endpgfonlayer
          \box\commtikz@savedfigbox
          \endgroup
        \else
          \commtikz@orig@tikz@finish
        \fi
      \else
        \commtikz@orig@tikz@finish
      \fi
    \fi
  }
\fi

% ==============================================================================
% 2. Signal Domain Line Styles, Waveguide Bends & Arrow Decorations
% ==============================================================================
\tikzset{
  every node/.append code={\commtikz@pathonbgfalse},
  behind nodes/.is if=commtikz@pathonbg,
  behind nodes/.default=true,
  % Base connection line style
  commtikz line/.style={
    behind nodes=true,
    line width=\pgfkeysvalueof{/tikz/commtikz/signal line width},
    line cap=butt,
    line join=miter,
  },
  commtikz arrow tip/.style={
    -{Latex[length=4.6pt, width=3.3pt]}
  },
  commtikz rev arrow tip/.style={
    {Latex[length=4.6pt, width=3.3pt]}-
  },
  % Signal domain styles
  opt/.style={draw=OPTcol, commtikz line},
  rf/.style={draw=RFcol, commtikz line},
  optrf/.style={draw=OPTRFcol, commtikz line},
  dc/.style={draw=DCcol, commtikz line},
  ctrl/.style={draw=CTRLcol, commtikz line, densely dashed},
  % Circular / rounded 90-deg and 180-deg bend helpers for ALL connection types
  bend/.style={rounded corners=#1},
  bend/.default=0.22cm,
  waveguide bend/.style={bend=#1},
  waveguide bend/.default=0.22cm,
  opt bend/.style={opt, bend=#1},
  opt bend/.default=0.22cm,
  rf bend/.style={rf, bend=#1},
  rf bend/.default=0.22cm,
  optrf bend/.style={optrf, bend=#1},
  optrf bend/.default=0.22cm,
  dc bend/.style={dc, bend=#1},
  dc bend/.default=0.22cm,
  ctrl bend/.style={ctrl, bend=#1},
  ctrl bend/.default=0.22cm,
  % Exact 180-degree semicircular U-turns (work on opt, rf, optrf, dc, ctrl, or any path)
  uturn/.style={
    to path={
      let \p1 = (\tikztostart),
          \p2 = (\tikztotarget),
          \n{r} = {0.5 * abs(\y2 - \y1)},
          \n{xm} = {max(\x1, \x2) + #1} in
      (\tikztostart) -- (\n{xm}, \y1)
      arc[start angle={(\y2 > \y1 ? -90 : 90)}, end angle={(\y2 > \y1 ? 90 : -90)}, radius=\n{r}]
      -- (\tikztotarget) \tikztonodes
    }
  },
  uturn/.default=0pt,
  uturn left/.style={
    to path={
      let \p1 = (\tikztostart),
          \p2 = (\tikztotarget),
          \n{r} = {0.5 * abs(\y2 - \y1)},
          \n{xm} = {min(\x1, \x2) - #1} in
      (\tikztostart) -- (\n{xm}, \y1)
      arc[start angle={(\y2 > \y1 ? 270 : 90)}, end angle={(\y2 > \y1 ? 90 : 270)}, radius=\n{r}]
      -- (\tikztotarget) \tikztonodes
    }
  },
  uturn left/.default=0pt,
  % Directed signal domain styles
  opt->/.style={opt, commtikz arrow tip},
  rf->/.style={rf, commtikz arrow tip},
  optrf->/.style={optrf, commtikz arrow tip},
  dc->/.style={dc, commtikz arrow tip},
  ctrl->/.style={ctrl, commtikz arrow tip},
  <-opt/.style={opt, commtikz rev arrow tip},
  <-rf/.style={rf, commtikz rev arrow tip},
  <-optrf/.style={optrf, commtikz rev arrow tip},
  % Mid-path arrow decorations (with automatic rounded-corner path resolution)
  commtikz round before decorate/.code={%
    \ifx\tikz@actions@path\@undefined\else
      \pgfprocessround{\tikz@actions@path}{\tikz@actions@path}%
    \fi
    \pgfsyssoftpath@getcurrentpath\commtikz@temppath
    \pgfprocessround{\commtikz@temppath}{\commtikz@temppath}%
    \pgfsyssoftpath@setcurrentpath\commtikz@temppath
    \pgfsetcornersarced{\pgfpointorigin}%
  },
  ->-/.style={
    decoration={
      markings,
      mark=at position #1 with {\pgfsetcornersarced{\pgfpointorigin}\arrow{Latex[length=4.6pt, width=3.3pt]}}
    },
    postaction={commtikz round before decorate, decorate}
  },
  ->-/.default=0.5,
  -<-/.style={
    decoration={
      markings,
      mark=at position #1 with {\pgfsetcornersarced{\pgfpointorigin}\arrowreversed{Latex[length=4.6pt, width=3.3pt]}}
    },
    postaction={commtikz round before decorate, decorate}
  },
  -<-/.default=0.5,
  ->>-/.style={
    decoration={
      markings,
      mark=at position 0.43 with {\pgfsetcornersarced{\pgfpointorigin}\arrow{Latex[length=4.6pt, width=3.3pt]}},
      mark=at position 0.57 with {\pgfsetcornersarced{\pgfpointorigin}\arrow{Latex[length=4.6pt, width=3.3pt]}}
    },
    postaction={commtikz round before decorate, decorate}
  },
  -<<-/.style={
    decoration={
      markings,
      mark=at position 0.43 with {\pgfsetcornersarced{\pgfpointorigin}\arrowreversed{Latex[length=4.6pt, width=3.3pt]}},
      mark=at position 0.57 with {\pgfsetcornersarced{\pgfpointorigin}\arrowreversed{Latex[length=4.6pt, width=3.3pt]}}
    },
    postaction={commtikz round before decorate, decorate}
  },
  % ==============================================================================
  % 3. Curved Connectors (arcz, arcn, arc2, arc3, arcloop)
  % ==============================================================================
  % arcz: Horizontal-to-horizontal S/Z curve
  arcz/.style={
    out=0,
    in=180,
    looseness=#1
  },
  arcz/.default=0.85,
  arcz rev/.style={
    out=180,
    in=0,
    looseness=#1
  },
  arcz rev/.default=0.85,
  % arcn: Vertical-to-vertical N curve
  arcn/.style={
    out=-90,
    in=90,
    looseness=#1
  },
  arcn/.default=0.85,
  arcn up/.style={
    out=90,
    in=-90,
    looseness=#1
  },
  arcn up/.default=0.85,
  % arc2: Symmetric C curve
  arc2/.style={
    bend left=#1
  },
  arc2/.default=30,
  arc2 right/.style={
    bend right=#1
  },
  arc2 right/.default=30,
  % arc3: Custom cubic curve (pass out=..., in=..., looseness=...)
  arc3/.style={
    out=45,
    in=135,
    #1
  },
  arc3/.default={},
  % arcloop: High-arched loop between two points
  arcloop/.style={
    out=75,
    in=105,
    looseness=#1
  },
  arcloop/.default=1.4,
  % Orthogonal and perpendicular routing helpers
  hv/.style={to path={-| (\tikztotarget) \tikztonodes}},
  vh/.style={to path={|- (\tikztotarget) \tikztonodes}},
  vperp/.style={to path={-- (\tikztostart |- \tikztotarget) \tikztonodes}},
  vperp to/.style={to path={(\tikztostart -| \tikztotarget) -- (\tikztotarget) \tikztonodes}},
  hperp/.style={to path={-- (\tikztostart -| \tikztotarget) \tikztonodes}},
  hperp to/.style={to path={(\tikztostart |- \tikztotarget) -- (\tikztotarget) \tikztonodes}},
  perp/.style={
    to path={
      let \p1 = (\tikztostart),
          \p2 = (\tikztotarget),
          \n{px} = {abs(\y2-\y1) >= abs(\x2-\x1) ? \x1 : \x2},
          \n{py} = {abs(\y2-\y1) >= abs(\x2-\x1) ? \y2 : \y1} in
      -- (\n{px}, \n{py}) \tikztonodes
    }
  },
  perp to/.style={
    to path={
      let \p1 = (\tikztostart),
          \p2 = (\tikztotarget),
          \n{px} = {abs(\y2-\y1) >= abs(\x2-\x1) ? \x2 : \x1},
          \n{py} = {abs(\y2-\y1) >= abs(\x2-\x1) ? \y1 : \y2} in
      (\n{px}, \n{py}) -- (\tikztotarget) \tikztonodes
    }
  },
  hvh/.style={
    to path={
      -- ($(\tikztostart)!#1!(\tikztotarget |- \tikztostart)$)
      |- (\tikztotarget) \tikztonodes
    }
  },
  hvh/.default=0.5,
  vhv/.style={
    to path={
      -- ($(\tikztostart)!#1!(\tikztostart |- \tikztotarget)$)
      -| (\tikztotarget) \tikztonodes
    }
  },
  vhv/.default=0.5,
  % ==============================================================================
  % 4. Inline Component Placement along Paths: to[inline={...}]
  % ==============================================================================
  commtikz save path color/.code={%
    % Resolve deferred color options in the path's scope, before node options
    % reset the stroke. This also handles color=... and color shorthands.
    \pgfinterruptpath
    \pgfscope
      \tikz@options
      % TikZ's bare color shorthands update the driver stroke state without
      % updating the pgfstrokecolor alias. Read the resolved driver color.
      \expandafter\let\csname\string\color@commtikz@resolved\endcsname\pgf@strokecolor@global
      \extractcolorspec{commtikz@resolved}{\commtikz@pathspec}%
      \edef\commtikz@next{\noexpand\endpgfscope\noexpand\endpgfinterruptpath
        \noexpand\definecolor{commtikz@pathcolor}\commtikz@pathspec}%
    \commtikz@next
    \def\commtikz@saved@pathcolor{commtikz@pathcolor}%
  },
  commtikz use path color/.code={%
    \ifx\commtikz@saved@pathcolor\pgfutil@empty
      \tikzset{draw=black}%
    \else
      \tikzset{draw=\commtikz@saved@pathcolor}%
    \fi
  },
  inline/.style={
    to path={
      \pgfextra{\tikzset{commtikz save path color}}%
      let \p{commtikz@d} = ($(\tikztotarget)-(\tikztostart)$),
          \n{commtikz@ang} = {atan2(\y{commtikz@d},\x{commtikz@d})} in
      node[at={($(\tikztostart)!0.5!(\tikztotarget)$)}, rotate=\n{commtikz@ang}, transform shape,
           solid, draw=black, text=black, name=commtikz@tmp@inline, #1, alias=commtikz@tmp@inline] {}
      (\tikztostart) -- (commtikz@tmp@inline.west)
      (commtikz@tmp@inline.east) -- (\tikztotarget)
      \tikztonodes
    }
  },
  inline upright/.style={
    to path={
      \pgfextra{\tikzset{commtikz save path color}}%
      node[at={($(\tikztostart)!0.5!(\tikztotarget)$)},
           solid, draw=black, text=black, name=commtikz@tmp@inline, #1, alias=commtikz@tmp@inline] {}
      (\tikztostart) -- (commtikz@tmp@inline.west)
      (commtikz@tmp@inline.east) -- (\tikztotarget)
      \tikztonodes
    }
  },
  % Convenient inline shorthands
  to coupler/.style={inline={coupler, commtikz use path color, #1}},
  to coupler/.default={},
  to couplercirc/.style={inline={couplercirc, commtikz use path color, #1}},
  to couplercirc/.default={},
  to fiber/.style={inline={fiber, commtikz use path color, #1}},
  to fiber/.default={},
  to fiberspool/.style={inline={fiberspool, commtikz use path color, #1}},
  to fiberspool/.default={},
  to polctrl/.style={inline={polctrl, #1}},
  to polctrl/.default={},
  to voa/.style={inline={voa, #1}},
  to voa/.default={},
  to fbg/.style={inline={fbg, #1}},
  to fbg/.default={},
  to isolator/.style={inline={isolator, #1}},
  to isolator/.default={},
  to pm/.style={inline={pm, #1}},
  to pm/.default={},
  to mzm/.style={inline={mzm, #1}},
  to mzm/.default={},
  to amp/.style={inline={amp, #1}},
  to amp/.default={},
  to edfa/.style={inline={edfa, #1}},
  to edfa/.default={},
  to filter/.style={inline={filter, #1}},
  to filter/.default={},
}
